\documentclass[conference]{IEEEtran}
\usepackage{amsmath,amssymb,booktabs,multirow,graphicx}
\usepackage[hidelinks]{hyperref}
\graphicspath{{../../figures/}}

\newcommand{\method}{FOWAC}
% RAW-V2 SAFETY: never replace these placeholders from legacy cumulative-mean logs.
\newcommand{\pending}{TBD}

\title{Structure-Referenced Open-World Few-Shot\\Class-Incremental Audio Learning}
\author{Anonymous ICASSP submission}

\begin{document}
\maketitle
\begin{abstract}
Open-world class-incremental audio learning must reject unknown inputs, discover their
latent categories, and learn them without erasing earlier knowledge. We study this joint
problem under mixed known/novel streams. A shared latent-structure reference bank replaces
the unsupported interpretation of late ResNet axes as physical time and frequency.
An exemplar-free uncertainty memory stores class moments and reliability, calibrates
five-shot covariance, and performs confidence-anchored hard statistical replay. On
LS-100, NS-100 and FSC-89, we compare recent continual category-discovery components
under a common four-session protocol with paired stream repeats. Corrected ten-stream
experiments report explicit round/session records; final aggregate values and paired
tests are inserted only after all raw-v2 runs complete. Reliability gates prevent
harmful replay when the base representation or variance estimate is inadequate.
\end{abstract}

\section{Introduction}
Few-shot class-incremental audio classification normally assumes labeled novel support,
whereas few-shot open-set recognition terminates after rejection. A deployed learner must
instead repeat detection, discovery and classifier expansion. This coupling exposes two
failures: rejected streams are contaminated by old classes, and noisy five-shot centers
increasingly compete with old prototypes.
Training-free calibration such as TEEN~\cite{wang2023teen} corrects biased few-shot
prototypes, while recent continual category-discovery methods address prediction/hardness
bias~\cite{ma2024happy}, projected distillation~\cite{rypesc2024camp}, online streams
~\cite{park2024ocgcd}, or covariance misalignment~\cite{dai2025vbcgcd}. The unified
open-set/continual formulation is also related to OpenIncrement~\cite{xu2023openincrement}
and the MetaGCD setting~\cite{wu2023metagcd}. Their visual protocols, however, do not jointly test
audio open-set rejection and five-shot category registration.

We make three contributions. First, we define a repeated mixed-stream audio protocol that
reports detection, discovery and continual classification jointly. Second, we learn one
Latent Structure Reference Bank (LSRB) during base training and reuse its stable
responses for later sessions; its centers are references in latent feature space, not
acoustic units. Third, we introduce uncertainty-aware moment memory and bounded hard
statistical replay. Our result ledger separates exact runs, official-code transfers,
component reimplementations and incompatible protocols.

\begin{figure*}[t]
\centering
\includegraphics[width=0.98\textwidth]{framework_overview_image2_v4.png}
\caption{Overview of \method{}. A dataset-level LSRB is learned once from latent
layer-4 descriptors. During each open-world session, registered samples bypass discovery,
whereas rejected candidates undergo capacity-constrained discovery and dual-prototype
expansion. Reliability-gated moment replay protects old prototypes while retaining
plasticity for new classes.}
\label{fig:framework}
\end{figure*}

\section{Method}
Let session zero contain labeled classes $\mathcal C_0$. Session $t$ presents an unlabeled
stream mixing registered classes and five novel classes. A ResNet-18 yields pooled semantic
embedding $z\in\mathbb R^{512}$ and layer-4 descriptors $\{f_i\}$. Normalized descriptors
from the base set are reservoir sampled and clustered once into $B=\{b_k\}_{k=1}^K$.
Their structural response is
\begin{equation}
a_{ik}=\frac{\exp(\cos(f_i,b_k)/\tau)}{\sum_j\exp(\cos(f_i,b_j)/\tau)},\qquad
h_k(x)=\frac{1}{HW}\sum_i a_{ik}.
\end{equation}
During base training, the bank is updated from cross-sample descriptors and its response
is fed back to the encoder through a structure-consistency term,
\begin{equation}
\mathcal L_{\mathrm{lsrb}}=\frac{1}{B}\sum_{k=1}^{B}
\left\|\operatorname{sg}(\bar h_k)-h_k(x)\right\|_2^2,
\end{equation}
where $\bar h_k$ is the batch reference response and $\operatorname{sg}$ stops
the target branch. The base objective is
$\mathcal L_{\mathrm{base}}=\mathcal L_{\mathrm{CE}}+
\lambda_s\mathcal L_{\mathrm{lsrb}}+\lambda_c\mathcal L_{\mathrm{center}}$.
After convergence, the bank is frozen only for cross-session evaluation, ensuring that
LSRB participates in representation learning rather than acting as offline post-processing.
The bank provides a common reference across samples and sessions. Novel candidates
are clustered in a normalized semantic--structural space. Candidate filtering can use a
fixed old-prototype similarity threshold or a stream-adaptive similarity quantile.

For each registered class, memory $M_c=(\mu_c,v_c,n_c,\rho_c)$ stores its normalized mean,
diagonal variance, count and reliability. A sparse novel estimate is shrunk toward related
old distributions:
\begin{equation}
\widetilde v_n=(1-\eta_n)\widehat v_n+eta_n\sum_c w_{nc}v_c,\quad
w_{nc}=\operatorname{softmax}(\cos(\widehat\mu_n,\mu_c)/\tau_m).
\end{equation}
Pseudo-features drawn from these moments are weighted by their maximum competing-prototype
similarity. Prototype optimization uses replay cross-entropy and a reliability anchor,
\begin{equation}
\mathcal L=\sum_i q_{y_i}\operatorname{CE}(\bar x_i,W)+
\lambda_a\sum_c(2-\rho_c)\lVert w_c-w_c^{old}\rVert_2^2,
\end{equation}
then writes back only a bounded residual
$W\leftarrow(1-\delta)W^{old}+\delta W^{rep}$. Gates disable replay when base accuracy is
below 0.7 or normalized variance is below $10^{-4}$.

Ordinary K-means can merge most sparse candidates into one cluster. We initialize with
K-means and alternate centroid updates with a global minimum-cost assignment
\begin{equation}
\min_A\sum_{ik}A_{ik}\lVert g_i-c_k\rVert^2,\quad
\sum_k A_{ik}=1,\quad \sum_i A_{ik}\in\{\lfloor N/K\rfloor,\lceil N/K\rceil\}.
\end{equation}
The Hungarian solver implements the assignment by expanding every center into capacity
slots. This uses the known $K$-way stream construction but no candidate labels.
In the released implementation, the historical flags ``DFSB'' and
``balanced\_kmeans/BCD'' refer to LSRB and CANA, respectively.

\section{Experiments}
LS-100 and NS-100 use 80 base classes; FSC-89 uses 69. Each has four 5-way 5-shot
incremental sessions. Every stream mixes five old and five novel classes with five samples
per class. Complete streams are repeated ten times with classifier reset, and comparisons
use paired streams. Audio is sampled at 16 kHz and converted to 128-bin log-Mel features
with 25-ms windows and 10-ms hops. Base optimization uses Nesterov SGD, learning rate
0.005, weight decay $5\times10^{-4}$ and batch size 128.

We report known/unknown recall, F1, AUROC, AUPR and FPR95; clustering ACC, NMI and ARI;
incremental-only, old-class and all-seen accuracy; final accuracy and forgetting. Baselines
include the project classifier, TEEN, VB-CGCD, OFCL, OPCR, YLOC and a Happy component
transfer. Image-oriented methods use the same frozen session features and are identified as
transfers rather than original-protocol reproductions.

\begin{table}[t]
\centering\small
\caption{Average incremental accuracy (\%). All main values use ten repeats.}
\label{tab:main}
\begin{tabular}{lccc}
\toprule Method & LS-100 & NS-100 & FSC-89\\
\midrule
TEEN & \pending & \pending & \pending\\
OFCL transfer & 50.81 & 59.86 & 0.73\\
OPCR transfer & 43.27 & 43.13 & 1.02\\
YLOC transfer & 52.17 & \textbf{60.17} & 0.00\\
Happy transfer & 37.46 & 28.34 & 5.56\\
\method{}-UMR & \pending & \pending & \pending\\
\method{}-DS & \pending & -- & --\\
\method{}-CANA & -- & -- & \pending\\
\bottomrule
\end{tabular}
\end{table}

Raw-v2 paired experiments for UMR, LSRB and CANA are in progress; means, sample standard
deviations, bootstrap intervals and paired tests will be generated directly from complete
round/session records. Unrestricted write-back and ungated
replay are negative ablations. Candidate diagnostics show that incremental accuracy is
most correlated with candidate purity on LS-100 and FSC-89 ($r=0.79$ and $0.66$), but
with novel-class coverage on NS-100 ($r=0.48$); this motivates dataset-independent,
stream-adaptive candidate control rather than a universal fixed threshold.
On FSC-89, CANA is evaluated with the same paired raw-v2 streams.
On NS-100 and FSC-89, reliability gates are designed to make UMR transparent
when replay statistics are inadequate; ungated replay is a negative ablation.

\begin{figure}[t]
\centering
\includegraphics[width=\columnwidth]{candidate_quality_correlation.pdf}
\caption{Candidate quality controls downstream incremental learning. Lines are least-squares
fits over repeated sessions; panels show the dominant significant correlate per dataset.}
\label{fig:candidate}
\end{figure}

\section{Conclusion}
Shared deep-feature structure supplies a technically defensible cross-session reference,
while uncertainty memory targets old/new prototype interference without retaining audio.
The reliability analysis also establishes when replay should be skipped rather than forced.

\bibliographystyle{IEEEtran}
\bibliography{references}
\end{document}
